% ============================================================
%  Theory Section — CALPHAD-Coupled Phase-Field Model
%  For: Cahn-Hilliard Spinodal Decomposition in
%       Multi-Component BCC High-Entropy Alloys
% ============================================================
%
% Usage: % ============================================================
%  Theory Section — CALPHAD-Coupled Phase-Field Model
%  For: Cahn-Hilliard Spinodal Decomposition in
%       Multi-Component BCC High-Entropy Alloys
% ============================================================
%
% Usage: % ============================================================
%  Theory Section — CALPHAD-Coupled Phase-Field Model
%  For: Cahn-Hilliard Spinodal Decomposition in
%       Multi-Component BCC High-Entropy Alloys
% ============================================================
%
% Usage: % ============================================================
%  Theory Section — CALPHAD-Coupled Phase-Field Model
%  For: Cahn-Hilliard Spinodal Decomposition in
%       Multi-Component BCC High-Entropy Alloys
% ============================================================
%
% Usage: \input{docs/theory_section.tex}  inside your article
% Packages required: amsmath, amssymb, bm, physics (or manual \dd)
% ============================================================

\section{Theory}
\label{sec:theory}

\subsection{Cahn-Hilliard Phase-Field Formulation}
\label{sec:ch_formulation}

We describe spinodal decomposition in an $N$-component alloy by a
set of $N-1$ conserved composition fields
$c_i(\bm{x},t) = x_i(\bm{x},t)$, $i = 1,\ldots,N-1$,
where $x_i$ is the mole fraction of element $i$.
The mole fraction of the reference element $N$ is recovered by the
closure relation $x_N = 1 - \sum_{i=1}^{N-1} c_i$.
The system evolves according to the Cahn-Hilliard equation
\begin{equation}
  \frac{\partial c_i}{\partial t}
  = \nabla \cdot \sum_{j=1}^{N-1}
      M_{ij}(\bm{c})\,\nabla\mu_j^{\mathrm{tot}},
  \qquad i = 1,\ldots,N-1,
  \label{eq:CH}
\end{equation}
where $\bm{M}(\bm{c})$ is the composition-dependent mobility tensor
and $\mu_j^{\mathrm{tot}}$ is the variational chemical potential of
field $j$ with respect to the free energy functional $\mathcal{F}$.

\subsubsection{Free Energy Functional}
The total free energy functional combines a bulk thermodynamic term
and an interfacial gradient penalty,
\begin{equation}
  \mathcal{F}[\bm{c}]
  = \int_\Omega \left[
      G_m(\bm{c})
      + \sum_{i=1}^{N-1} \frac{\kappa_i}{2}\,|\nabla c_i|^2
    \right] d^3x,
  \label{eq:free_energy}
\end{equation}
where $G_m(\bm{c})$ is the molar Gibbs energy of the alloy (in
J~mol$^{-1}$) and $\kappa_i$ (J~m$^2$~mol$^{-1}$) is the gradient
energy coefficient for field $i$.  The variational derivative of
Eq.~\eqref{eq:free_energy} gives
\begin{equation}
  \mu_j^{\mathrm{tot}} = \frac{\delta \mathcal{F}}{\delta c_j}
  = \underbrace{\frac{\partial G_m}{\partial c_j}}_{\mu_j}
    - \kappa_j\,\nabla^2 c_j,
  \label{eq:mu_var}
\end{equation}
where the \emph{bulk chemical potential}
$\mu_j = \partial G_m/\partial c_j$
is the thermodynamic driving force for diffusion and the
$-\kappa_j\nabla^2 c_j$ term penalises sharp interfaces.

\subsubsection{Degenerate Mobility}
A degenerate mobility tensor is employed to enforce the physical
constraint that diffusion vanishes in the limit of a pure component,
\begin{equation}
  M_{ij}(\bm{c}) = \mathring{M}_{ij}
    \prod_{k=1}^{N} x_k(\bm{c}),
  \label{eq:mobility}
\end{equation}
where $\mathring{M}_{ij}$ (m$^5$~J$^{-1}$~s$^{-1}$) is the
constant (bare) mobility matrix supplied as a model parameter.
The product over all $N$ mole fractions, including the reference
element, vanishes whenever any $x_k \to 0$.  In practice, the
product is replaced by its spatial average over the domain,
$\langle \prod_k x_k \rangle$, to avoid an explicit spatial
coupling and maintain the pseudo-spectral efficiency of the solver.

\subsubsection{Pseudo-Spectral Semi-Implicit Time Integration}
Equation~\eqref{eq:CH} is solved in Fourier space on a uniform
$N_g\times N_g$ grid with domain size $L_x$.  Diagonal mobility
terms are treated implicitly (unconditionally damping high-wavenumber
modes) while off-diagonal cross-terms are treated explicitly:
\begin{equation}
  \hat{c}_i^{\,n+1}
  = \frac{
      \hat{c}_i^{\,n}
      + \Delta t\,k^2 \sum_j M_{ij}^{\mathrm{eff}}\,\hat{\mu}_{j,\mathrm{bulk}}^n
      - \Delta t\,k^4 \sum_{j\neq i} M_{ij}^{\mathrm{eff}}\,\kappa_j\,\hat{c}_j^{\,n}
    }{
      1 + \Delta t\,k^4\,M_{ii}^{\mathrm{eff}}\,\kappa_i
    },
  \label{eq:semi_implicit}
\end{equation}
where $\hat{(\cdot)}$ denotes the Fourier transform,
$k^2 = k_x^2 + k_y^2$ is the squared wavenumber magnitude, and
$M_{ij}^{\mathrm{eff}} = \mathring{M}_{ij}\,\langle\prod_k x_k\rangle$
is the spatially averaged effective mobility.
The stable time step is estimated as
\begin{equation}
  \Delta t_{\mathrm{stable}}
  = \frac{\alpha\,\Delta x^2}
         {\max_i(\mathring{M}_{ii})\,|\lambda_{\max}|},
  \label{eq:dt_stable}
\end{equation}
where $\alpha < 1$ is a safety factor, $\Delta x = L_x/N_g$ is the
grid spacing, and $|\lambda_{\max}|$ is the largest absolute
eigenvalue of the Hessian matrix $H_{ij} = \partial^2 G_m/\partial
c_i\,\partial c_j$ evaluated in a neighbourhood of the initial
mean composition $\bm{c}_0$.

% -----------------------------------------------------------
\subsection{CALPHAD-Derived Thermodynamic Driving Forces}
\label{sec:calphad}

The bulk chemical potential $\mu_j = \partial G_m/\partial c_j$ in
Eq.~\eqref{eq:mu_var} must be evaluated at every spatial grid point
at every time step.  This evaluation is the dominant computational
cost and its accuracy determines whether the simulated microstructural
evolution is physically consistent with the underlying thermodynamics.
Direct calls to a CALPHAD solver (pycalphad~\cite{pycalphad}) at each
point and step are prohibitively expensive.  We therefore construct
smooth surrogate functions from thermodynamic data pre-computed on a
dense composition grid.

\subsubsection{Pre-computation Strategy}
For each alloy system and temperature of interest, the following
quantities are evaluated on a simplex grid
$\{x_j^{(k)}\}$ spanning the physically relevant composition space,
using pycalphad's \texttt{IsolatedPhase} framework (which constrains
the evaluation to the single-phase BCC$_{\rm A2}$ free energy
surface):
\begin{itemize}
  \item $G_m^{(k)}$: molar Gibbs energy (J~mol$^{-1}$);
  \item $\mu_i^{(k)} = \texttt{MU}(i)$: absolute chemical potential of
        element $i$ (J~mol$^{-1}$);
  \item $\Lambda_{ij}^{(k)} = \texttt{MU}(i)\texttt{.X}(j)
        \equiv \partial\mu_i/\partial x_j$:
        Jansson second derivative~\cite{jansson1984} (J~mol$^{-1}$).
\end{itemize}
The computed data are stored in HDF5 format and loaded once at the
start of a simulation run.  All subsequent thermodynamic queries
during time integration are answered by interpolation.

\subsubsection{Surrogate Construction}
The surrogate is built as a single callable object
(\texttt{GmInterpolant}) that maps an arbitrary spatial composition
field $\bm{c}(\bm{x})$ to $G_m$, $\mu_j$, and $H_{ij}$.  Three
backend interpolation schemes are supported, selected on a
per-simulation basis:

\begin{itemize}
  \item \textbf{Binary ($N=2$, one independent field):}
        CubicSpline (scipy)~\cite{scipy} on the 1-D grid of
        $x_{\rm Nb}$ values.
  \item \textbf{Ternary ($N=3$, two independent fields):}
        CloughTocher C1 triangular interpolant~\cite{clough_tocher}
        on the scattered simplex data, with a nearest-neighbour
        fallback for the rare points that fall outside the convex
        hull of the training set.
  \item \textbf{Quaternary and beyond ($N \geq 4$):}
        Two-stage scheme: (i)~a thin-plate-spline
        RBFInterpolator~\cite{rbf_tps} is fitted once on the
        scattered CALPHAD data; (ii)~it is evaluated on a dense
        regular grid of $n_g^{N-1}$ points; (iii)~a
        RegularGridInterpolator (trilinear) is constructed on the
        resulting grid for $O(\log n_g)$ lookup during simulation.
\end{itemize}
The same surrogate class is used for $G_m$, the individual $\mu_i$,
and the Hessian components $H_{ij}$, so all three derivative methods
described below share identical interpolation backends.

% -----------------------------------------------------------
\subsection{Derivative Methods for Bulk Chemical Potentials}
\label{sec:deriv_methods}

The phase-field model requires two thermodynamic quantities at each
grid point per time step: the bulk driving force
$\mu_j = \partial G_m/\partial c_j$ (first derivative) and the
Hessian $H_{ij} = \partial^2 G_m/\partial c_i\partial c_j$ (second
derivative, needed for stability analysis and time-step selection).
Three strategies are implemented for obtaining these quantities,
each representing a different trade-off between accuracy, data
requirements, and computational overhead.

\subsubsection{Method~I: Finite Differences on the $G_m$ Surrogate (\texttt{fd})}
\label{sec:fd}

The simplest approach is to differentiate the $G_m$ surrogate
numerically.  Central finite differences give
\begin{align}
  \mu_j^{\rm fd}
  &= \frac{G_m(\bm{c}+\varepsilon\bm{e}_j)
           - G_m(\bm{c}-\varepsilon\bm{e}_j)}{2\varepsilon}
  + O(\varepsilon^2),
  \label{eq:fd_mu}
  \\
  H_{ij}^{\rm fd}
  &= \begin{cases}
       \dfrac{G_m(\bm{c}+\varepsilon\bm{e}_i)
              - 2G_m(\bm{c})
              + G_m(\bm{c}-\varepsilon\bm{e}_i)}{\varepsilon^2}
       + O(\varepsilon^2) & i = j,\\[8pt]
       \dfrac{G_m(\bm{c}+\varepsilon\bm{e}_i+\varepsilon\bm{e}_j)
              - G_m(\bm{c}+\varepsilon\bm{e}_i-\varepsilon\bm{e}_j)
              - G_m(\bm{c}-\varepsilon\bm{e}_i+\varepsilon\bm{e}_j)
              + G_m(\bm{c}-\varepsilon\bm{e}_i-\varepsilon\bm{e}_j)}{4\varepsilon^2}
       + O(\varepsilon^2) & i \neq j,
     \end{cases}
  \label{eq:fd_H}
\end{align}
where $\bm{e}_j$ is the unit vector in the $j$-th composition
direction and $\varepsilon = 10^{-4}$ mol fraction.
This method requires only $G_m$ data (no chemical potential storage)
but introduces truncation errors that compound with the smoothing
bias of the $G_m$ surrogate.  In particular, the second-order
accurate stencil in Eq.~\eqref{eq:fd_H} can accumulate errors of
order $\varepsilon^2 G_m'''/G_m$ near the simplex boundaries where
$G_m$ is steep; our validation on the Nb--V binary at $500\,^\circ$C
yields an RMSE of $442\,\text{kJ\,mol}^{-1}$ in $H$ relative to the
directly computed CALPHAD Hessian (Section~\ref{sec:validation}).
This error is dominated by the entropy divergence
$H \to +\infty$ as $x \to 0$ or $x \to 1$ and its propagation
into the FD stencil.

\subsubsection{Method~II: Direct Interpolation of Jansson Derivatives (\texttt{jansson})}
\label{sec:jansson}

Under the simplex constraint $x_N = 1 - \sum_{i=1}^{N-1} c_i$, the
bulk chemical potential of field $j$ that enters the Cahn-Hilliard
equation is the \emph{relative} chemical potential
\begin{equation}
  \mu_j = \frac{\partial G_m}{\partial c_j}\bigg|_{\rm simplex}
        = \bar{\mu}_j - \bar{\mu}_N,
  \label{eq:mu_rel}
\end{equation}
where $\bar{\mu}_i = \texttt{MU}(i)$ is the absolute chemical
potential of element $i$~\cite{deFontaine1979}.  This identity
follows directly from the Gibbs-Duhem equation and is exact within
the single-phase thermodynamic model.

Similarly, the Hessian element is
\begin{equation}
  H_{ij}
  = \frac{\partial^2 G_m}{\partial c_i \partial c_j}\bigg|_{\rm simplex}
  = \frac{\partial(\bar{\mu}_i - \bar{\mu}_N)}{\partial x_j}
  = \Lambda_{ij} - \Lambda_{Nj},
  \label{eq:H_jansson}
\end{equation}
where $\Lambda_{ij} = \partial\bar{\mu}_i/\partial x_j$ is the
Jansson second derivative~\cite{jansson1984} computed by pycalphad
and stored in the HDF5 table (Section~\ref{sec:calphad}).

In the \texttt{jansson} mode, separate smooth surrogate functions are
built for each $\mu_j$ (Eq.~\ref{eq:mu_rel}) and each $H_{ij}$
(Eq.~\ref{eq:H_jansson}) using the same backend as the $G_m$
interpolant.  During time integration, all thermodynamic quantities
are obtained from these surrogates by simple function evaluation with
no numerical differentiation.  Because the Jansson derivatives are
computed by pycalphad using the analytic tangent to the Gibbs energy
surface, this method is free of FD truncation error and satisfies the
Gibbs-Duhem relation exactly at every composition point.  Our
validation confirms RMSE $= 0$ for both $\mu$ and $H$ relative to
the raw CALPHAD data (Section~\ref{sec:validation}).

For an $N$-component system, the \texttt{jansson} mode builds
$(N-1)$ surrogates for $\mu_j$ and $(N-1)^2$ surrogates for
$H_{ij}$, for a total of $(N-1)N$ additional interpolants
beyond the $G_m$ surface.  For quaternary alloys ($N=4$) this
amounts to 3 + 9 = 12 additional RBF$\to$RGI surfaces, each
constructed once at simulation startup.

\subsubsection{Method~III: Hybrid — Jansson $\mu$, Finite-Difference $H$ (\texttt{hybrid})}
\label{sec:hybrid}

The \texttt{hybrid} mode combines the best attributes of the two
preceding methods.  The bulk driving force $\mu_j$ is obtained by
direct interpolation of the stored Jansson $\bar{\mu}$ data
(Eq.~\ref{eq:mu_rel}), which eliminates the first-derivative
truncation error from the time-evolution equation.  The Hessian
$H_{ij}$ is obtained by finite differences on the $G_m$ surrogate
(Eq.~\ref{eq:fd_H}), which is used only for stability analysis and
time-step estimation rather than for the per-step driving force.
This is appropriate when the \texttt{dMU} dataset is absent or when
the raw Jansson data is suspected to be noisy (e.g., near a
miscibility gap or phase boundary where the single-phase free energy
surface becomes nearly flat and pycalphad's numerical derivatives are
less reliable).

A practical summary of the three methods is given in
Table~\ref{tab:methods}.

\begin{table}[htbp]
  \centering
  \caption{%
    Summary of derivative methods for bulk chemical potentials and
    Hessian in the CALPHAD-coupled Cahn-Hilliard model.
    RMSE values are for the Nb--V binary BCC$_{\rm A2}$ system at
    $T = 500\,^\circ$C with $n = 200$ composition grid points.
  }
  \label{tab:methods}
  \setlength{\tabcolsep}{6pt}
  \renewcommand{\arraystretch}{1.25}
  \begin{tabular}{llllcc}
    \hline
    Mode & $\mu_j$ source & $H_{ij}$ source
         & HDF5 data needed
         & RMSE$_\mu$ & RMSE$_H$ \\
    \hline
    \texttt{fd}
      & FD on $G_m$ (Eq.~\ref{eq:fd_mu})
      & FD on $G_m$ (Eq.~\ref{eq:fd_H})
      & \texttt{/GM}
      & 302 & 442\,200 \\
    \texttt{jansson}
      & Eq.~\ref{eq:mu_rel}
      & Eq.~\ref{eq:H_jansson}
      & \texttt{/GM}, \texttt{/MU}, \texttt{/dMU}
      & \textbf{0} & \textbf{0} \\
    \texttt{hybrid}
      & Eq.~\ref{eq:mu_rel}
      & FD on $G_m$ (Eq.~\ref{eq:fd_H})
      & \texttt{/GM}, \texttt{/MU}
      & \textbf{0} & 442\,200 \\
    \hline
  \end{tabular}
  \par\smallskip
  \footnotesize RMSE in J~mol$^{-1}$; \textbf{bold} = exact match
  with CALPHAD data.
\end{table}

% -----------------------------------------------------------
\subsection{Validation of Derivative Methods}
\label{sec:validation}

Figure~\ref{fig:deriv_methods} compares the three modes against raw
CALPHAD data for the Nb--V BCC$_{\rm A2}$ system at
$T = 500\,^\circ$C as a function of the vanadium mole fraction
$x_{\rm V}$.  The molar Gibbs energy $G_m$ (panel~a) is identical
across all modes because it is always obtained from the same
$G_m$ interpolant.  The relative chemical potential
$\mu_{\rm Nb} - \mu_{\rm V}$ (panel~b) shows that \texttt{fd}
deviates visibly from the CALPHAD reference near
$x_{\rm V} \approx 0.2$--$0.3$, the composition range where the
$G_m$ spline has the steepest first derivative.  Both \texttt{jansson}
and \texttt{hybrid} reproduce the CALPHAD chemical potentials exactly
at all tabulated compositions.

The most significant differences appear in the Hessian
$\partial^2 G_m/\partial x_{\rm Nb}^2$ (panel~c).  The y-axis is
clipped to $\pm 150\,\text{kJ\,mol}^{-1}$ to expose the spinodal
region (shaded pink, $H < 0$); the ideal-entropy divergence
$H \to +\infty$ as $x_{\rm V} \to 0$ or $1$ is annotated at both
edges.  The \texttt{fd} curve deviates substantially from the
CALPHAD reference due to the compounding of spline smoothing bias
and FD truncation error, with an RMSE of $442\,\text{kJ\,mol}^{-1}$
(Table~\ref{tab:methods}).  Since the Hessian controls the spinodal
decomposition wavelength via $\lambda^* = 4\pi\sqrt{\kappa/|H|}$,
errors of this magnitude translate directly into inaccurate
predictions of the dominant microstructural length scale.  The
\texttt{jansson} mode eliminates this error entirely, at the cost of
pre-computing the Jansson derivative dataset (step~1 of the workflow,
which is performed once and stored on disk).

We therefore recommend \texttt{jansson} mode as the default for
production runs, and \texttt{hybrid} mode as a fall-back when
the \texttt{dMU} dataset is unavailable or when Jansson data quality
near a phase boundary is suspect.  The \texttt{fd} mode is retained
for backward compatibility and for rapid prototyping where
thermodynamic accuracy is secondary to workflow simplicity.

% -----------------------------------------------------------
\subsection{Spinodal Stability and Interface Length Scales}
\label{sec:length_scales}

The spinodal decomposition instability is determined by the minimum
eigenvalue $\lambda_{\min}$ of the Hessian $\bm{H}(\bm{c}_0)$
evaluated at the mean composition $\bm{c}_0$.  If
$\lambda_{\min} < 0$, the alloy is spinodally unstable and phase
separation proceeds spontaneously.  The dominant decomposition
wavelength and interface width are
\begin{equation}
  \lambda_i^* = 4\pi\sqrt{\frac{\kappa_i}{|\lambda_{\min}|}},
  \qquad
  \delta_i = \sqrt{\frac{\kappa_i}{|\lambda_{\min}|}},
  \qquad
  i = 1,\ldots,N-1,
  \label{eq:length_scales}
\end{equation}
and the shortest unstable wavelength (critical wavelength below which
all modes are stable) is
\begin{equation}
  \lambda_{c,i} = 2\pi\sqrt{\frac{2\kappa_i}{|\lambda_{\min}|}}.
  \label{eq:lambda_c}
\end{equation}
Proper resolution of the interface requires $\delta_i > 5\Delta x$;
the gradient energy coefficient $\kappa_i$ must be chosen
accordingly.  A numerical stability criterion for the semi-implicit
scheme (Eq.~\ref{eq:semi_implicit}) is
$M_{ii}|\lambda_{\min}|\Delta t/\Delta x^2 < 0.5$, which is
evaluated per-field and reported alongside $\Delta t_{\rm stable}$
(Eq.~\ref{eq:dt_stable}) in the simulation parameter summary written
to disk.
  inside your article
% Packages required: amsmath, amssymb, bm, physics (or manual \dd)
% ============================================================

\section{Theory}
\label{sec:theory}

\subsection{Cahn-Hilliard Phase-Field Formulation}
\label{sec:ch_formulation}

We describe spinodal decomposition in an $N$-component alloy by a
set of $N-1$ conserved composition fields
$c_i(\bm{x},t) = x_i(\bm{x},t)$, $i = 1,\ldots,N-1$,
where $x_i$ is the mole fraction of element $i$.
The mole fraction of the reference element $N$ is recovered by the
closure relation $x_N = 1 - \sum_{i=1}^{N-1} c_i$.
The system evolves according to the Cahn-Hilliard equation
\begin{equation}
  \frac{\partial c_i}{\partial t}
  = \nabla \cdot \sum_{j=1}^{N-1}
      M_{ij}(\bm{c})\,\nabla\mu_j^{\mathrm{tot}},
  \qquad i = 1,\ldots,N-1,
  \label{eq:CH}
\end{equation}
where $\bm{M}(\bm{c})$ is the composition-dependent mobility tensor
and $\mu_j^{\mathrm{tot}}$ is the variational chemical potential of
field $j$ with respect to the free energy functional $\mathcal{F}$.

\subsubsection{Free Energy Functional}
The total free energy functional combines a bulk thermodynamic term
and an interfacial gradient penalty,
\begin{equation}
  \mathcal{F}[\bm{c}]
  = \int_\Omega \left[
      G_m(\bm{c})
      + \sum_{i=1}^{N-1} \frac{\kappa_i}{2}\,|\nabla c_i|^2
    \right] d^3x,
  \label{eq:free_energy}
\end{equation}
where $G_m(\bm{c})$ is the molar Gibbs energy of the alloy (in
J~mol$^{-1}$) and $\kappa_i$ (J~m$^2$~mol$^{-1}$) is the gradient
energy coefficient for field $i$.  The variational derivative of
Eq.~\eqref{eq:free_energy} gives
\begin{equation}
  \mu_j^{\mathrm{tot}} = \frac{\delta \mathcal{F}}{\delta c_j}
  = \underbrace{\frac{\partial G_m}{\partial c_j}}_{\mu_j}
    - \kappa_j\,\nabla^2 c_j,
  \label{eq:mu_var}
\end{equation}
where the \emph{bulk chemical potential}
$\mu_j = \partial G_m/\partial c_j$
is the thermodynamic driving force for diffusion and the
$-\kappa_j\nabla^2 c_j$ term penalises sharp interfaces.

\subsubsection{Degenerate Mobility}
A degenerate mobility tensor is employed to enforce the physical
constraint that diffusion vanishes in the limit of a pure component,
\begin{equation}
  M_{ij}(\bm{c}) = \mathring{M}_{ij}
    \prod_{k=1}^{N} x_k(\bm{c}),
  \label{eq:mobility}
\end{equation}
where $\mathring{M}_{ij}$ (m$^5$~J$^{-1}$~s$^{-1}$) is the
constant (bare) mobility matrix supplied as a model parameter.
The product over all $N$ mole fractions, including the reference
element, vanishes whenever any $x_k \to 0$.  In practice, the
product is replaced by its spatial average over the domain,
$\langle \prod_k x_k \rangle$, to avoid an explicit spatial
coupling and maintain the pseudo-spectral efficiency of the solver.

\subsubsection{Pseudo-Spectral Semi-Implicit Time Integration}
Equation~\eqref{eq:CH} is solved in Fourier space on a uniform
$N_g\times N_g$ grid with domain size $L_x$.  Diagonal mobility
terms are treated implicitly (unconditionally damping high-wavenumber
modes) while off-diagonal cross-terms are treated explicitly:
\begin{equation}
  \hat{c}_i^{\,n+1}
  = \frac{
      \hat{c}_i^{\,n}
      + \Delta t\,k^2 \sum_j M_{ij}^{\mathrm{eff}}\,\hat{\mu}_{j,\mathrm{bulk}}^n
      - \Delta t\,k^4 \sum_{j\neq i} M_{ij}^{\mathrm{eff}}\,\kappa_j\,\hat{c}_j^{\,n}
    }{
      1 + \Delta t\,k^4\,M_{ii}^{\mathrm{eff}}\,\kappa_i
    },
  \label{eq:semi_implicit}
\end{equation}
where $\hat{(\cdot)}$ denotes the Fourier transform,
$k^2 = k_x^2 + k_y^2$ is the squared wavenumber magnitude, and
$M_{ij}^{\mathrm{eff}} = \mathring{M}_{ij}\,\langle\prod_k x_k\rangle$
is the spatially averaged effective mobility.
The stable time step is estimated as
\begin{equation}
  \Delta t_{\mathrm{stable}}
  = \frac{\alpha\,\Delta x^2}
         {\max_i(\mathring{M}_{ii})\,|\lambda_{\max}|},
  \label{eq:dt_stable}
\end{equation}
where $\alpha < 1$ is a safety factor, $\Delta x = L_x/N_g$ is the
grid spacing, and $|\lambda_{\max}|$ is the largest absolute
eigenvalue of the Hessian matrix $H_{ij} = \partial^2 G_m/\partial
c_i\,\partial c_j$ evaluated in a neighbourhood of the initial
mean composition $\bm{c}_0$.

% -----------------------------------------------------------
\subsection{CALPHAD-Derived Thermodynamic Driving Forces}
\label{sec:calphad}

The bulk chemical potential $\mu_j = \partial G_m/\partial c_j$ in
Eq.~\eqref{eq:mu_var} must be evaluated at every spatial grid point
at every time step.  This evaluation is the dominant computational
cost and its accuracy determines whether the simulated microstructural
evolution is physically consistent with the underlying thermodynamics.
Direct calls to a CALPHAD solver (pycalphad~\cite{pycalphad}) at each
point and step are prohibitively expensive.  We therefore construct
smooth surrogate functions from thermodynamic data pre-computed on a
dense composition grid.

\subsubsection{Pre-computation Strategy}
For each alloy system and temperature of interest, the following
quantities are evaluated on a simplex grid
$\{x_j^{(k)}\}$ spanning the physically relevant composition space,
using pycalphad's \texttt{IsolatedPhase} framework (which constrains
the evaluation to the single-phase BCC$_{\rm A2}$ free energy
surface):
\begin{itemize}
  \item $G_m^{(k)}$: molar Gibbs energy (J~mol$^{-1}$);
  \item $\mu_i^{(k)} = \texttt{MU}(i)$: absolute chemical potential of
        element $i$ (J~mol$^{-1}$);
  \item $\Lambda_{ij}^{(k)} = \texttt{MU}(i)\texttt{.X}(j)
        \equiv \partial\mu_i/\partial x_j$:
        Jansson second derivative~\cite{jansson1984} (J~mol$^{-1}$).
\end{itemize}
The computed data are stored in HDF5 format and loaded once at the
start of a simulation run.  All subsequent thermodynamic queries
during time integration are answered by interpolation.

\subsubsection{Surrogate Construction}
The surrogate is built as a single callable object
(\texttt{GmInterpolant}) that maps an arbitrary spatial composition
field $\bm{c}(\bm{x})$ to $G_m$, $\mu_j$, and $H_{ij}$.  Three
backend interpolation schemes are supported, selected on a
per-simulation basis:

\begin{itemize}
  \item \textbf{Binary ($N=2$, one independent field):}
        CubicSpline (scipy)~\cite{scipy} on the 1-D grid of
        $x_{\rm Nb}$ values.
  \item \textbf{Ternary ($N=3$, two independent fields):}
        CloughTocher C1 triangular interpolant~\cite{clough_tocher}
        on the scattered simplex data, with a nearest-neighbour
        fallback for the rare points that fall outside the convex
        hull of the training set.
  \item \textbf{Quaternary and beyond ($N \geq 4$):}
        Two-stage scheme: (i)~a thin-plate-spline
        RBFInterpolator~\cite{rbf_tps} is fitted once on the
        scattered CALPHAD data; (ii)~it is evaluated on a dense
        regular grid of $n_g^{N-1}$ points; (iii)~a
        RegularGridInterpolator (trilinear) is constructed on the
        resulting grid for $O(\log n_g)$ lookup during simulation.
\end{itemize}
The same surrogate class is used for $G_m$, the individual $\mu_i$,
and the Hessian components $H_{ij}$, so all three derivative methods
described below share identical interpolation backends.

% -----------------------------------------------------------
\subsection{Derivative Methods for Bulk Chemical Potentials}
\label{sec:deriv_methods}

The phase-field model requires two thermodynamic quantities at each
grid point per time step: the bulk driving force
$\mu_j = \partial G_m/\partial c_j$ (first derivative) and the
Hessian $H_{ij} = \partial^2 G_m/\partial c_i\partial c_j$ (second
derivative, needed for stability analysis and time-step selection).
Three strategies are implemented for obtaining these quantities,
each representing a different trade-off between accuracy, data
requirements, and computational overhead.

\subsubsection{Method~I: Finite Differences on the $G_m$ Surrogate (\texttt{fd})}
\label{sec:fd}

The simplest approach is to differentiate the $G_m$ surrogate
numerically.  Central finite differences give
\begin{align}
  \mu_j^{\rm fd}
  &= \frac{G_m(\bm{c}+\varepsilon\bm{e}_j)
           - G_m(\bm{c}-\varepsilon\bm{e}_j)}{2\varepsilon}
  + O(\varepsilon^2),
  \label{eq:fd_mu}
  \\
  H_{ij}^{\rm fd}
  &= \begin{cases}
       \dfrac{G_m(\bm{c}+\varepsilon\bm{e}_i)
              - 2G_m(\bm{c})
              + G_m(\bm{c}-\varepsilon\bm{e}_i)}{\varepsilon^2}
       + O(\varepsilon^2) & i = j,\\[8pt]
       \dfrac{G_m(\bm{c}+\varepsilon\bm{e}_i+\varepsilon\bm{e}_j)
              - G_m(\bm{c}+\varepsilon\bm{e}_i-\varepsilon\bm{e}_j)
              - G_m(\bm{c}-\varepsilon\bm{e}_i+\varepsilon\bm{e}_j)
              + G_m(\bm{c}-\varepsilon\bm{e}_i-\varepsilon\bm{e}_j)}{4\varepsilon^2}
       + O(\varepsilon^2) & i \neq j,
     \end{cases}
  \label{eq:fd_H}
\end{align}
where $\bm{e}_j$ is the unit vector in the $j$-th composition
direction and $\varepsilon = 10^{-4}$ mol fraction.
This method requires only $G_m$ data (no chemical potential storage)
but introduces truncation errors that compound with the smoothing
bias of the $G_m$ surrogate.  In particular, the second-order
accurate stencil in Eq.~\eqref{eq:fd_H} can accumulate errors of
order $\varepsilon^2 G_m'''/G_m$ near the simplex boundaries where
$G_m$ is steep; our validation on the Nb--V binary at $500\,^\circ$C
yields an RMSE of $442\,\text{kJ\,mol}^{-1}$ in $H$ relative to the
directly computed CALPHAD Hessian (Section~\ref{sec:validation}).
This error is dominated by the entropy divergence
$H \to +\infty$ as $x \to 0$ or $x \to 1$ and its propagation
into the FD stencil.

\subsubsection{Method~II: Direct Interpolation of Jansson Derivatives (\texttt{jansson})}
\label{sec:jansson}

Under the simplex constraint $x_N = 1 - \sum_{i=1}^{N-1} c_i$, the
bulk chemical potential of field $j$ that enters the Cahn-Hilliard
equation is the \emph{relative} chemical potential
\begin{equation}
  \mu_j = \frac{\partial G_m}{\partial c_j}\bigg|_{\rm simplex}
        = \bar{\mu}_j - \bar{\mu}_N,
  \label{eq:mu_rel}
\end{equation}
where $\bar{\mu}_i = \texttt{MU}(i)$ is the absolute chemical
potential of element $i$~\cite{deFontaine1979}.  This identity
follows directly from the Gibbs-Duhem equation and is exact within
the single-phase thermodynamic model.

Similarly, the Hessian element is
\begin{equation}
  H_{ij}
  = \frac{\partial^2 G_m}{\partial c_i \partial c_j}\bigg|_{\rm simplex}
  = \frac{\partial(\bar{\mu}_i - \bar{\mu}_N)}{\partial x_j}
  = \Lambda_{ij} - \Lambda_{Nj},
  \label{eq:H_jansson}
\end{equation}
where $\Lambda_{ij} = \partial\bar{\mu}_i/\partial x_j$ is the
Jansson second derivative~\cite{jansson1984} computed by pycalphad
and stored in the HDF5 table (Section~\ref{sec:calphad}).

In the \texttt{jansson} mode, separate smooth surrogate functions are
built for each $\mu_j$ (Eq.~\ref{eq:mu_rel}) and each $H_{ij}$
(Eq.~\ref{eq:H_jansson}) using the same backend as the $G_m$
interpolant.  During time integration, all thermodynamic quantities
are obtained from these surrogates by simple function evaluation with
no numerical differentiation.  Because the Jansson derivatives are
computed by pycalphad using the analytic tangent to the Gibbs energy
surface, this method is free of FD truncation error and satisfies the
Gibbs-Duhem relation exactly at every composition point.  Our
validation confirms RMSE $= 0$ for both $\mu$ and $H$ relative to
the raw CALPHAD data (Section~\ref{sec:validation}).

For an $N$-component system, the \texttt{jansson} mode builds
$(N-1)$ surrogates for $\mu_j$ and $(N-1)^2$ surrogates for
$H_{ij}$, for a total of $(N-1)N$ additional interpolants
beyond the $G_m$ surface.  For quaternary alloys ($N=4$) this
amounts to 3 + 9 = 12 additional RBF$\to$RGI surfaces, each
constructed once at simulation startup.

\subsubsection{Method~III: Hybrid — Jansson $\mu$, Finite-Difference $H$ (\texttt{hybrid})}
\label{sec:hybrid}

The \texttt{hybrid} mode combines the best attributes of the two
preceding methods.  The bulk driving force $\mu_j$ is obtained by
direct interpolation of the stored Jansson $\bar{\mu}$ data
(Eq.~\ref{eq:mu_rel}), which eliminates the first-derivative
truncation error from the time-evolution equation.  The Hessian
$H_{ij}$ is obtained by finite differences on the $G_m$ surrogate
(Eq.~\ref{eq:fd_H}), which is used only for stability analysis and
time-step estimation rather than for the per-step driving force.
This is appropriate when the \texttt{dMU} dataset is absent or when
the raw Jansson data is suspected to be noisy (e.g., near a
miscibility gap or phase boundary where the single-phase free energy
surface becomes nearly flat and pycalphad's numerical derivatives are
less reliable).

A practical summary of the three methods is given in
Table~\ref{tab:methods}.

\begin{table}[htbp]
  \centering
  \caption{%
    Summary of derivative methods for bulk chemical potentials and
    Hessian in the CALPHAD-coupled Cahn-Hilliard model.
    RMSE values are for the Nb--V binary BCC$_{\rm A2}$ system at
    $T = 500\,^\circ$C with $n = 200$ composition grid points.
  }
  \label{tab:methods}
  \setlength{\tabcolsep}{6pt}
  \renewcommand{\arraystretch}{1.25}
  \begin{tabular}{llllcc}
    \hline
    Mode & $\mu_j$ source & $H_{ij}$ source
         & HDF5 data needed
         & RMSE$_\mu$ & RMSE$_H$ \\
    \hline
    \texttt{fd}
      & FD on $G_m$ (Eq.~\ref{eq:fd_mu})
      & FD on $G_m$ (Eq.~\ref{eq:fd_H})
      & \texttt{/GM}
      & 302 & 442\,200 \\
    \texttt{jansson}
      & Eq.~\ref{eq:mu_rel}
      & Eq.~\ref{eq:H_jansson}
      & \texttt{/GM}, \texttt{/MU}, \texttt{/dMU}
      & \textbf{0} & \textbf{0} \\
    \texttt{hybrid}
      & Eq.~\ref{eq:mu_rel}
      & FD on $G_m$ (Eq.~\ref{eq:fd_H})
      & \texttt{/GM}, \texttt{/MU}
      & \textbf{0} & 442\,200 \\
    \hline
  \end{tabular}
  \par\smallskip
  \footnotesize RMSE in J~mol$^{-1}$; \textbf{bold} = exact match
  with CALPHAD data.
\end{table}

% -----------------------------------------------------------
\subsection{Validation of Derivative Methods}
\label{sec:validation}

Figure~\ref{fig:deriv_methods} compares the three modes against raw
CALPHAD data for the Nb--V BCC$_{\rm A2}$ system at
$T = 500\,^\circ$C as a function of the vanadium mole fraction
$x_{\rm V}$.  The molar Gibbs energy $G_m$ (panel~a) is identical
across all modes because it is always obtained from the same
$G_m$ interpolant.  The relative chemical potential
$\mu_{\rm Nb} - \mu_{\rm V}$ (panel~b) shows that \texttt{fd}
deviates visibly from the CALPHAD reference near
$x_{\rm V} \approx 0.2$--$0.3$, the composition range where the
$G_m$ spline has the steepest first derivative.  Both \texttt{jansson}
and \texttt{hybrid} reproduce the CALPHAD chemical potentials exactly
at all tabulated compositions.

The most significant differences appear in the Hessian
$\partial^2 G_m/\partial x_{\rm Nb}^2$ (panel~c).  The y-axis is
clipped to $\pm 150\,\text{kJ\,mol}^{-1}$ to expose the spinodal
region (shaded pink, $H < 0$); the ideal-entropy divergence
$H \to +\infty$ as $x_{\rm V} \to 0$ or $1$ is annotated at both
edges.  The \texttt{fd} curve deviates substantially from the
CALPHAD reference due to the compounding of spline smoothing bias
and FD truncation error, with an RMSE of $442\,\text{kJ\,mol}^{-1}$
(Table~\ref{tab:methods}).  Since the Hessian controls the spinodal
decomposition wavelength via $\lambda^* = 4\pi\sqrt{\kappa/|H|}$,
errors of this magnitude translate directly into inaccurate
predictions of the dominant microstructural length scale.  The
\texttt{jansson} mode eliminates this error entirely, at the cost of
pre-computing the Jansson derivative dataset (step~1 of the workflow,
which is performed once and stored on disk).

We therefore recommend \texttt{jansson} mode as the default for
production runs, and \texttt{hybrid} mode as a fall-back when
the \texttt{dMU} dataset is unavailable or when Jansson data quality
near a phase boundary is suspect.  The \texttt{fd} mode is retained
for backward compatibility and for rapid prototyping where
thermodynamic accuracy is secondary to workflow simplicity.

% -----------------------------------------------------------
\subsection{Spinodal Stability and Interface Length Scales}
\label{sec:length_scales}

The spinodal decomposition instability is determined by the minimum
eigenvalue $\lambda_{\min}$ of the Hessian $\bm{H}(\bm{c}_0)$
evaluated at the mean composition $\bm{c}_0$.  If
$\lambda_{\min} < 0$, the alloy is spinodally unstable and phase
separation proceeds spontaneously.  The dominant decomposition
wavelength and interface width are
\begin{equation}
  \lambda_i^* = 4\pi\sqrt{\frac{\kappa_i}{|\lambda_{\min}|}},
  \qquad
  \delta_i = \sqrt{\frac{\kappa_i}{|\lambda_{\min}|}},
  \qquad
  i = 1,\ldots,N-1,
  \label{eq:length_scales}
\end{equation}
and the shortest unstable wavelength (critical wavelength below which
all modes are stable) is
\begin{equation}
  \lambda_{c,i} = 2\pi\sqrt{\frac{2\kappa_i}{|\lambda_{\min}|}}.
  \label{eq:lambda_c}
\end{equation}
Proper resolution of the interface requires $\delta_i > 5\Delta x$;
the gradient energy coefficient $\kappa_i$ must be chosen
accordingly.  A numerical stability criterion for the semi-implicit
scheme (Eq.~\ref{eq:semi_implicit}) is
$M_{ii}|\lambda_{\min}|\Delta t/\Delta x^2 < 0.5$, which is
evaluated per-field and reported alongside $\Delta t_{\rm stable}$
(Eq.~\ref{eq:dt_stable}) in the simulation parameter summary written
to disk.
  inside your article
% Packages required: amsmath, amssymb, bm, physics (or manual \dd)
% ============================================================

\section{Theory}
\label{sec:theory}

\subsection{Cahn-Hilliard Phase-Field Formulation}
\label{sec:ch_formulation}

We describe spinodal decomposition in an $N$-component alloy by a
set of $N-1$ conserved composition fields
$c_i(\bm{x},t) = x_i(\bm{x},t)$, $i = 1,\ldots,N-1$,
where $x_i$ is the mole fraction of element $i$.
The mole fraction of the reference element $N$ is recovered by the
closure relation $x_N = 1 - \sum_{i=1}^{N-1} c_i$.
The system evolves according to the Cahn-Hilliard equation
\begin{equation}
  \frac{\partial c_i}{\partial t}
  = \nabla \cdot \sum_{j=1}^{N-1}
      M_{ij}(\bm{c})\,\nabla\mu_j^{\mathrm{tot}},
  \qquad i = 1,\ldots,N-1,
  \label{eq:CH}
\end{equation}
where $\bm{M}(\bm{c})$ is the composition-dependent mobility tensor
and $\mu_j^{\mathrm{tot}}$ is the variational chemical potential of
field $j$ with respect to the free energy functional $\mathcal{F}$.

\subsubsection{Free Energy Functional}
The total free energy functional combines a bulk thermodynamic term
and an interfacial gradient penalty,
\begin{equation}
  \mathcal{F}[\bm{c}]
  = \int_\Omega \left[
      G_m(\bm{c})
      + \sum_{i=1}^{N-1} \frac{\kappa_i}{2}\,|\nabla c_i|^2
    \right] d^3x,
  \label{eq:free_energy}
\end{equation}
where $G_m(\bm{c})$ is the molar Gibbs energy of the alloy (in
J~mol$^{-1}$) and $\kappa_i$ (J~m$^2$~mol$^{-1}$) is the gradient
energy coefficient for field $i$.  The variational derivative of
Eq.~\eqref{eq:free_energy} gives
\begin{equation}
  \mu_j^{\mathrm{tot}} = \frac{\delta \mathcal{F}}{\delta c_j}
  = \underbrace{\frac{\partial G_m}{\partial c_j}}_{\mu_j}
    - \kappa_j\,\nabla^2 c_j,
  \label{eq:mu_var}
\end{equation}
where the \emph{bulk chemical potential}
$\mu_j = \partial G_m/\partial c_j$
is the thermodynamic driving force for diffusion and the
$-\kappa_j\nabla^2 c_j$ term penalises sharp interfaces.

\subsubsection{Degenerate Mobility}
A degenerate mobility tensor is employed to enforce the physical
constraint that diffusion vanishes in the limit of a pure component,
\begin{equation}
  M_{ij}(\bm{c}) = \mathring{M}_{ij}
    \prod_{k=1}^{N} x_k(\bm{c}),
  \label{eq:mobility}
\end{equation}
where $\mathring{M}_{ij}$ (m$^5$~J$^{-1}$~s$^{-1}$) is the
constant (bare) mobility matrix supplied as a model parameter.
The product over all $N$ mole fractions, including the reference
element, vanishes whenever any $x_k \to 0$.  In practice, the
product is replaced by its spatial average over the domain,
$\langle \prod_k x_k \rangle$, to avoid an explicit spatial
coupling and maintain the pseudo-spectral efficiency of the solver.

\subsubsection{Pseudo-Spectral Semi-Implicit Time Integration}
Equation~\eqref{eq:CH} is solved in Fourier space on a uniform
$N_g\times N_g$ grid with domain size $L_x$.  Diagonal mobility
terms are treated implicitly (unconditionally damping high-wavenumber
modes) while off-diagonal cross-terms are treated explicitly:
\begin{equation}
  \hat{c}_i^{\,n+1}
  = \frac{
      \hat{c}_i^{\,n}
      + \Delta t\,k^2 \sum_j M_{ij}^{\mathrm{eff}}\,\hat{\mu}_{j,\mathrm{bulk}}^n
      - \Delta t\,k^4 \sum_{j\neq i} M_{ij}^{\mathrm{eff}}\,\kappa_j\,\hat{c}_j^{\,n}
    }{
      1 + \Delta t\,k^4\,M_{ii}^{\mathrm{eff}}\,\kappa_i
    },
  \label{eq:semi_implicit}
\end{equation}
where $\hat{(\cdot)}$ denotes the Fourier transform,
$k^2 = k_x^2 + k_y^2$ is the squared wavenumber magnitude, and
$M_{ij}^{\mathrm{eff}} = \mathring{M}_{ij}\,\langle\prod_k x_k\rangle$
is the spatially averaged effective mobility.
The stable time step is estimated as
\begin{equation}
  \Delta t_{\mathrm{stable}}
  = \frac{\alpha\,\Delta x^2}
         {\max_i(\mathring{M}_{ii})\,|\lambda_{\max}|},
  \label{eq:dt_stable}
\end{equation}
where $\alpha < 1$ is a safety factor, $\Delta x = L_x/N_g$ is the
grid spacing, and $|\lambda_{\max}|$ is the largest absolute
eigenvalue of the Hessian matrix $H_{ij} = \partial^2 G_m/\partial
c_i\,\partial c_j$ evaluated in a neighbourhood of the initial
mean composition $\bm{c}_0$.

% -----------------------------------------------------------
\subsection{CALPHAD-Derived Thermodynamic Driving Forces}
\label{sec:calphad}

The bulk chemical potential $\mu_j = \partial G_m/\partial c_j$ in
Eq.~\eqref{eq:mu_var} must be evaluated at every spatial grid point
at every time step.  This evaluation is the dominant computational
cost and its accuracy determines whether the simulated microstructural
evolution is physically consistent with the underlying thermodynamics.
Direct calls to a CALPHAD solver (pycalphad~\cite{pycalphad}) at each
point and step are prohibitively expensive.  We therefore construct
smooth surrogate functions from thermodynamic data pre-computed on a
dense composition grid.

\subsubsection{Pre-computation Strategy}
For each alloy system and temperature of interest, the following
quantities are evaluated on a simplex grid
$\{x_j^{(k)}\}$ spanning the physically relevant composition space,
using pycalphad's \texttt{IsolatedPhase} framework (which constrains
the evaluation to the single-phase BCC$_{\rm A2}$ free energy
surface):
\begin{itemize}
  \item $G_m^{(k)}$: molar Gibbs energy (J~mol$^{-1}$);
  \item $\mu_i^{(k)} = \texttt{MU}(i)$: absolute chemical potential of
        element $i$ (J~mol$^{-1}$);
  \item $\Lambda_{ij}^{(k)} = \texttt{MU}(i)\texttt{.X}(j)
        \equiv \partial\mu_i/\partial x_j$:
        Jansson second derivative~\cite{jansson1984} (J~mol$^{-1}$).
\end{itemize}
The computed data are stored in HDF5 format and loaded once at the
start of a simulation run.  All subsequent thermodynamic queries
during time integration are answered by interpolation.

\subsubsection{Surrogate Construction}
The surrogate is built as a single callable object
(\texttt{GmInterpolant}) that maps an arbitrary spatial composition
field $\bm{c}(\bm{x})$ to $G_m$, $\mu_j$, and $H_{ij}$.  Three
backend interpolation schemes are supported, selected on a
per-simulation basis:

\begin{itemize}
  \item \textbf{Binary ($N=2$, one independent field):}
        CubicSpline (scipy)~\cite{scipy} on the 1-D grid of
        $x_{\rm Nb}$ values.
  \item \textbf{Ternary ($N=3$, two independent fields):}
        CloughTocher C1 triangular interpolant~\cite{clough_tocher}
        on the scattered simplex data, with a nearest-neighbour
        fallback for the rare points that fall outside the convex
        hull of the training set.
  \item \textbf{Quaternary and beyond ($N \geq 4$):}
        Two-stage scheme: (i)~a thin-plate-spline
        RBFInterpolator~\cite{rbf_tps} is fitted once on the
        scattered CALPHAD data; (ii)~it is evaluated on a dense
        regular grid of $n_g^{N-1}$ points; (iii)~a
        RegularGridInterpolator (trilinear) is constructed on the
        resulting grid for $O(\log n_g)$ lookup during simulation.
\end{itemize}
The same surrogate class is used for $G_m$, the individual $\mu_i$,
and the Hessian components $H_{ij}$, so all three derivative methods
described below share identical interpolation backends.

% -----------------------------------------------------------
\subsection{Derivative Methods for Bulk Chemical Potentials}
\label{sec:deriv_methods}

The phase-field model requires two thermodynamic quantities at each
grid point per time step: the bulk driving force
$\mu_j = \partial G_m/\partial c_j$ (first derivative) and the
Hessian $H_{ij} = \partial^2 G_m/\partial c_i\partial c_j$ (second
derivative, needed for stability analysis and time-step selection).
Three strategies are implemented for obtaining these quantities,
each representing a different trade-off between accuracy, data
requirements, and computational overhead.

\subsubsection{Method~I: Finite Differences on the $G_m$ Surrogate (\texttt{fd})}
\label{sec:fd}

The simplest approach is to differentiate the $G_m$ surrogate
numerically.  Central finite differences give
\begin{align}
  \mu_j^{\rm fd}
  &= \frac{G_m(\bm{c}+\varepsilon\bm{e}_j)
           - G_m(\bm{c}-\varepsilon\bm{e}_j)}{2\varepsilon}
  + O(\varepsilon^2),
  \label{eq:fd_mu}
  \\
  H_{ij}^{\rm fd}
  &= \begin{cases}
       \dfrac{G_m(\bm{c}+\varepsilon\bm{e}_i)
              - 2G_m(\bm{c})
              + G_m(\bm{c}-\varepsilon\bm{e}_i)}{\varepsilon^2}
       + O(\varepsilon^2) & i = j,\\[8pt]
       \dfrac{G_m(\bm{c}+\varepsilon\bm{e}_i+\varepsilon\bm{e}_j)
              - G_m(\bm{c}+\varepsilon\bm{e}_i-\varepsilon\bm{e}_j)
              - G_m(\bm{c}-\varepsilon\bm{e}_i+\varepsilon\bm{e}_j)
              + G_m(\bm{c}-\varepsilon\bm{e}_i-\varepsilon\bm{e}_j)}{4\varepsilon^2}
       + O(\varepsilon^2) & i \neq j,
     \end{cases}
  \label{eq:fd_H}
\end{align}
where $\bm{e}_j$ is the unit vector in the $j$-th composition
direction and $\varepsilon = 10^{-4}$ mol fraction.
This method requires only $G_m$ data (no chemical potential storage)
but introduces truncation errors that compound with the smoothing
bias of the $G_m$ surrogate.  In particular, the second-order
accurate stencil in Eq.~\eqref{eq:fd_H} can accumulate errors of
order $\varepsilon^2 G_m'''/G_m$ near the simplex boundaries where
$G_m$ is steep; our validation on the Nb--V binary at $500\,^\circ$C
yields an RMSE of $442\,\text{kJ\,mol}^{-1}$ in $H$ relative to the
directly computed CALPHAD Hessian (Section~\ref{sec:validation}).
This error is dominated by the entropy divergence
$H \to +\infty$ as $x \to 0$ or $x \to 1$ and its propagation
into the FD stencil.

\subsubsection{Method~II: Direct Interpolation of Jansson Derivatives (\texttt{jansson})}
\label{sec:jansson}

Under the simplex constraint $x_N = 1 - \sum_{i=1}^{N-1} c_i$, the
bulk chemical potential of field $j$ that enters the Cahn-Hilliard
equation is the \emph{relative} chemical potential
\begin{equation}
  \mu_j = \frac{\partial G_m}{\partial c_j}\bigg|_{\rm simplex}
        = \bar{\mu}_j - \bar{\mu}_N,
  \label{eq:mu_rel}
\end{equation}
where $\bar{\mu}_i = \texttt{MU}(i)$ is the absolute chemical
potential of element $i$~\cite{deFontaine1979}.  This identity
follows directly from the Gibbs-Duhem equation and is exact within
the single-phase thermodynamic model.

Similarly, the Hessian element is
\begin{equation}
  H_{ij}
  = \frac{\partial^2 G_m}{\partial c_i \partial c_j}\bigg|_{\rm simplex}
  = \frac{\partial(\bar{\mu}_i - \bar{\mu}_N)}{\partial x_j}
  = \Lambda_{ij} - \Lambda_{Nj},
  \label{eq:H_jansson}
\end{equation}
where $\Lambda_{ij} = \partial\bar{\mu}_i/\partial x_j$ is the
Jansson second derivative~\cite{jansson1984} computed by pycalphad
and stored in the HDF5 table (Section~\ref{sec:calphad}).

In the \texttt{jansson} mode, separate smooth surrogate functions are
built for each $\mu_j$ (Eq.~\ref{eq:mu_rel}) and each $H_{ij}$
(Eq.~\ref{eq:H_jansson}) using the same backend as the $G_m$
interpolant.  During time integration, all thermodynamic quantities
are obtained from these surrogates by simple function evaluation with
no numerical differentiation.  Because the Jansson derivatives are
computed by pycalphad using the analytic tangent to the Gibbs energy
surface, this method is free of FD truncation error and satisfies the
Gibbs-Duhem relation exactly at every composition point.  Our
validation confirms RMSE $= 0$ for both $\mu$ and $H$ relative to
the raw CALPHAD data (Section~\ref{sec:validation}).

For an $N$-component system, the \texttt{jansson} mode builds
$(N-1)$ surrogates for $\mu_j$ and $(N-1)^2$ surrogates for
$H_{ij}$, for a total of $(N-1)N$ additional interpolants
beyond the $G_m$ surface.  For quaternary alloys ($N=4$) this
amounts to 3 + 9 = 12 additional RBF$\to$RGI surfaces, each
constructed once at simulation startup.

\subsubsection{Method~III: Hybrid — Jansson $\mu$, Finite-Difference $H$ (\texttt{hybrid})}
\label{sec:hybrid}

The \texttt{hybrid} mode combines the best attributes of the two
preceding methods.  The bulk driving force $\mu_j$ is obtained by
direct interpolation of the stored Jansson $\bar{\mu}$ data
(Eq.~\ref{eq:mu_rel}), which eliminates the first-derivative
truncation error from the time-evolution equation.  The Hessian
$H_{ij}$ is obtained by finite differences on the $G_m$ surrogate
(Eq.~\ref{eq:fd_H}), which is used only for stability analysis and
time-step estimation rather than for the per-step driving force.
This is appropriate when the \texttt{dMU} dataset is absent or when
the raw Jansson data is suspected to be noisy (e.g., near a
miscibility gap or phase boundary where the single-phase free energy
surface becomes nearly flat and pycalphad's numerical derivatives are
less reliable).

A practical summary of the three methods is given in
Table~\ref{tab:methods}.

\begin{table}[htbp]
  \centering
  \caption{%
    Summary of derivative methods for bulk chemical potentials and
    Hessian in the CALPHAD-coupled Cahn-Hilliard model.
    RMSE values are for the Nb--V binary BCC$_{\rm A2}$ system at
    $T = 500\,^\circ$C with $n = 200$ composition grid points.
  }
  \label{tab:methods}
  \setlength{\tabcolsep}{6pt}
  \renewcommand{\arraystretch}{1.25}
  \begin{tabular}{llllcc}
    \hline
    Mode & $\mu_j$ source & $H_{ij}$ source
         & HDF5 data needed
         & RMSE$_\mu$ & RMSE$_H$ \\
    \hline
    \texttt{fd}
      & FD on $G_m$ (Eq.~\ref{eq:fd_mu})
      & FD on $G_m$ (Eq.~\ref{eq:fd_H})
      & \texttt{/GM}
      & 302 & 442\,200 \\
    \texttt{jansson}
      & Eq.~\ref{eq:mu_rel}
      & Eq.~\ref{eq:H_jansson}
      & \texttt{/GM}, \texttt{/MU}, \texttt{/dMU}
      & \textbf{0} & \textbf{0} \\
    \texttt{hybrid}
      & Eq.~\ref{eq:mu_rel}
      & FD on $G_m$ (Eq.~\ref{eq:fd_H})
      & \texttt{/GM}, \texttt{/MU}
      & \textbf{0} & 442\,200 \\
    \hline
  \end{tabular}
  \par\smallskip
  \footnotesize RMSE in J~mol$^{-1}$; \textbf{bold} = exact match
  with CALPHAD data.
\end{table}

% -----------------------------------------------------------
\subsection{Validation of Derivative Methods}
\label{sec:validation}

Figure~\ref{fig:deriv_methods} compares the three modes against raw
CALPHAD data for the Nb--V BCC$_{\rm A2}$ system at
$T = 500\,^\circ$C as a function of the vanadium mole fraction
$x_{\rm V}$.  The molar Gibbs energy $G_m$ (panel~a) is identical
across all modes because it is always obtained from the same
$G_m$ interpolant.  The relative chemical potential
$\mu_{\rm Nb} - \mu_{\rm V}$ (panel~b) shows that \texttt{fd}
deviates visibly from the CALPHAD reference near
$x_{\rm V} \approx 0.2$--$0.3$, the composition range where the
$G_m$ spline has the steepest first derivative.  Both \texttt{jansson}
and \texttt{hybrid} reproduce the CALPHAD chemical potentials exactly
at all tabulated compositions.

The most significant differences appear in the Hessian
$\partial^2 G_m/\partial x_{\rm Nb}^2$ (panel~c).  The y-axis is
clipped to $\pm 150\,\text{kJ\,mol}^{-1}$ to expose the spinodal
region (shaded pink, $H < 0$); the ideal-entropy divergence
$H \to +\infty$ as $x_{\rm V} \to 0$ or $1$ is annotated at both
edges.  The \texttt{fd} curve deviates substantially from the
CALPHAD reference due to the compounding of spline smoothing bias
and FD truncation error, with an RMSE of $442\,\text{kJ\,mol}^{-1}$
(Table~\ref{tab:methods}).  Since the Hessian controls the spinodal
decomposition wavelength via $\lambda^* = 4\pi\sqrt{\kappa/|H|}$,
errors of this magnitude translate directly into inaccurate
predictions of the dominant microstructural length scale.  The
\texttt{jansson} mode eliminates this error entirely, at the cost of
pre-computing the Jansson derivative dataset (step~1 of the workflow,
which is performed once and stored on disk).

We therefore recommend \texttt{jansson} mode as the default for
production runs, and \texttt{hybrid} mode as a fall-back when
the \texttt{dMU} dataset is unavailable or when Jansson data quality
near a phase boundary is suspect.  The \texttt{fd} mode is retained
for backward compatibility and for rapid prototyping where
thermodynamic accuracy is secondary to workflow simplicity.

% -----------------------------------------------------------
\subsection{Spinodal Stability and Interface Length Scales}
\label{sec:length_scales}

The spinodal decomposition instability is determined by the minimum
eigenvalue $\lambda_{\min}$ of the Hessian $\bm{H}(\bm{c}_0)$
evaluated at the mean composition $\bm{c}_0$.  If
$\lambda_{\min} < 0$, the alloy is spinodally unstable and phase
separation proceeds spontaneously.  The dominant decomposition
wavelength and interface width are
\begin{equation}
  \lambda_i^* = 4\pi\sqrt{\frac{\kappa_i}{|\lambda_{\min}|}},
  \qquad
  \delta_i = \sqrt{\frac{\kappa_i}{|\lambda_{\min}|}},
  \qquad
  i = 1,\ldots,N-1,
  \label{eq:length_scales}
\end{equation}
and the shortest unstable wavelength (critical wavelength below which
all modes are stable) is
\begin{equation}
  \lambda_{c,i} = 2\pi\sqrt{\frac{2\kappa_i}{|\lambda_{\min}|}}.
  \label{eq:lambda_c}
\end{equation}
Proper resolution of the interface requires $\delta_i > 5\Delta x$;
the gradient energy coefficient $\kappa_i$ must be chosen
accordingly.  A numerical stability criterion for the semi-implicit
scheme (Eq.~\ref{eq:semi_implicit}) is
$M_{ii}|\lambda_{\min}|\Delta t/\Delta x^2 < 0.5$, which is
evaluated per-field and reported alongside $\Delta t_{\rm stable}$
(Eq.~\ref{eq:dt_stable}) in the simulation parameter summary written
to disk.
  inside your article
% Packages required: amsmath, amssymb, bm, physics (or manual \dd)
% ============================================================

\section{Theory}
\label{sec:theory}

\subsection{Cahn-Hilliard Phase-Field Formulation}
\label{sec:ch_formulation}

We describe spinodal decomposition in an $N$-component alloy by a
set of $N-1$ conserved composition fields
$c_i(\bm{x},t) = x_i(\bm{x},t)$, $i = 1,\ldots,N-1$,
where $x_i$ is the mole fraction of element $i$.
The mole fraction of the reference element $N$ is recovered by the
closure relation $x_N = 1 - \sum_{i=1}^{N-1} c_i$.
The system evolves according to the Cahn-Hilliard equation
\begin{equation}
  \frac{\partial c_i}{\partial t}
  = \nabla \cdot \sum_{j=1}^{N-1}
      M_{ij}(\bm{c})\,\nabla\mu_j^{\mathrm{tot}},
  \qquad i = 1,\ldots,N-1,
  \label{eq:CH}
\end{equation}
where $\bm{M}(\bm{c})$ is the composition-dependent mobility tensor
and $\mu_j^{\mathrm{tot}}$ is the variational chemical potential of
field $j$ with respect to the free energy functional $\mathcal{F}$.

\subsubsection{Free Energy Functional}
The total free energy functional combines a bulk thermodynamic term
and an interfacial gradient penalty,
\begin{equation}
  \mathcal{F}[\bm{c}]
  = \int_\Omega \left[
      G_m(\bm{c})
      + \sum_{i=1}^{N-1} \frac{\kappa_i}{2}\,|\nabla c_i|^2
    \right] d^3x,
  \label{eq:free_energy}
\end{equation}
where $G_m(\bm{c})$ is the molar Gibbs energy of the alloy (in
J~mol$^{-1}$) and $\kappa_i$ (J~m$^2$~mol$^{-1}$) is the gradient
energy coefficient for field $i$.  The variational derivative of
Eq.~\eqref{eq:free_energy} gives
\begin{equation}
  \mu_j^{\mathrm{tot}} = \frac{\delta \mathcal{F}}{\delta c_j}
  = \underbrace{\frac{\partial G_m}{\partial c_j}}_{\mu_j}
    - \kappa_j\,\nabla^2 c_j,
  \label{eq:mu_var}
\end{equation}
where the \emph{bulk chemical potential}
$\mu_j = \partial G_m/\partial c_j$
is the thermodynamic driving force for diffusion and the
$-\kappa_j\nabla^2 c_j$ term penalises sharp interfaces.

\subsubsection{Degenerate Mobility}
A degenerate mobility tensor is employed to enforce the physical
constraint that diffusion vanishes in the limit of a pure component,
\begin{equation}
  M_{ij}(\bm{c}) = \mathring{M}_{ij}
    \prod_{k=1}^{N} x_k(\bm{c}),
  \label{eq:mobility}
\end{equation}
where $\mathring{M}_{ij}$ (m$^5$~J$^{-1}$~s$^{-1}$) is the
constant (bare) mobility matrix supplied as a model parameter.
The product over all $N$ mole fractions, including the reference
element, vanishes whenever any $x_k \to 0$.  In practice, the
product is replaced by its spatial average over the domain,
$\langle \prod_k x_k \rangle$, to avoid an explicit spatial
coupling and maintain the pseudo-spectral efficiency of the solver.

\subsubsection{Pseudo-Spectral Semi-Implicit Time Integration}
Equation~\eqref{eq:CH} is solved in Fourier space on a uniform
$N_g\times N_g$ grid with domain size $L_x$.  Diagonal mobility
terms are treated implicitly (unconditionally damping high-wavenumber
modes) while off-diagonal cross-terms are treated explicitly:
\begin{equation}
  \hat{c}_i^{\,n+1}
  = \frac{
      \hat{c}_i^{\,n}
      + \Delta t\,k^2 \sum_j M_{ij}^{\mathrm{eff}}\,\hat{\mu}_{j,\mathrm{bulk}}^n
      - \Delta t\,k^4 \sum_{j\neq i} M_{ij}^{\mathrm{eff}}\,\kappa_j\,\hat{c}_j^{\,n}
    }{
      1 + \Delta t\,k^4\,M_{ii}^{\mathrm{eff}}\,\kappa_i
    },
  \label{eq:semi_implicit}
\end{equation}
where $\hat{(\cdot)}$ denotes the Fourier transform,
$k^2 = k_x^2 + k_y^2$ is the squared wavenumber magnitude, and
$M_{ij}^{\mathrm{eff}} = \mathring{M}_{ij}\,\langle\prod_k x_k\rangle$
is the spatially averaged effective mobility.
The stable time step is estimated as
\begin{equation}
  \Delta t_{\mathrm{stable}}
  = \frac{\alpha\,\Delta x^2}
         {\max_i(\mathring{M}_{ii})\,|\lambda_{\max}|},
  \label{eq:dt_stable}
\end{equation}
where $\alpha < 1$ is a safety factor, $\Delta x = L_x/N_g$ is the
grid spacing, and $|\lambda_{\max}|$ is the largest absolute
eigenvalue of the Hessian matrix $H_{ij} = \partial^2 G_m/\partial
c_i\,\partial c_j$ evaluated in a neighbourhood of the initial
mean composition $\bm{c}_0$.

% -----------------------------------------------------------
\subsection{CALPHAD-Derived Thermodynamic Driving Forces}
\label{sec:calphad}

The bulk chemical potential $\mu_j = \partial G_m/\partial c_j$ in
Eq.~\eqref{eq:mu_var} must be evaluated at every spatial grid point
at every time step.  This evaluation is the dominant computational
cost and its accuracy determines whether the simulated microstructural
evolution is physically consistent with the underlying thermodynamics.
Direct calls to a CALPHAD solver (pycalphad~\cite{pycalphad}) at each
point and step are prohibitively expensive.  We therefore construct
smooth surrogate functions from thermodynamic data pre-computed on a
dense composition grid.

\subsubsection{Pre-computation Strategy}
For each alloy system and temperature of interest, the following
quantities are evaluated on a simplex grid
$\{x_j^{(k)}\}$ spanning the physically relevant composition space,
using pycalphad's \texttt{IsolatedPhase} framework (which constrains
the evaluation to the single-phase BCC$_{\rm A2}$ free energy
surface):
\begin{itemize}
  \item $G_m^{(k)}$: molar Gibbs energy (J~mol$^{-1}$);
  \item $\mu_i^{(k)} = \texttt{MU}(i)$: absolute chemical potential of
        element $i$ (J~mol$^{-1}$);
  \item $\Lambda_{ij}^{(k)} = \texttt{MU}(i)\texttt{.X}(j)
        \equiv \partial\mu_i/\partial x_j$:
        Jansson second derivative~\cite{jansson1984} (J~mol$^{-1}$).
\end{itemize}
The computed data are stored in HDF5 format and loaded once at the
start of a simulation run.  All subsequent thermodynamic queries
during time integration are answered by interpolation.

\subsubsection{Surrogate Construction}
The surrogate is built as a single callable object
(\texttt{GmInterpolant}) that maps an arbitrary spatial composition
field $\bm{c}(\bm{x})$ to $G_m$, $\mu_j$, and $H_{ij}$.  Three
backend interpolation schemes are supported, selected on a
per-simulation basis:

\begin{itemize}
  \item \textbf{Binary ($N=2$, one independent field):}
        CubicSpline (scipy)~\cite{scipy} on the 1-D grid of
        $x_{\rm Nb}$ values.
  \item \textbf{Ternary ($N=3$, two independent fields):}
        CloughTocher C1 triangular interpolant~\cite{clough_tocher}
        on the scattered simplex data, with a nearest-neighbour
        fallback for the rare points that fall outside the convex
        hull of the training set.
  \item \textbf{Quaternary and beyond ($N \geq 4$):}
        Two-stage scheme: (i)~a thin-plate-spline
        RBFInterpolator~\cite{rbf_tps} is fitted once on the
        scattered CALPHAD data; (ii)~it is evaluated on a dense
        regular grid of $n_g^{N-1}$ points; (iii)~a
        RegularGridInterpolator (trilinear) is constructed on the
        resulting grid for $O(\log n_g)$ lookup during simulation.
\end{itemize}
The same surrogate class is used for $G_m$, the individual $\mu_i$,
and the Hessian components $H_{ij}$, so all three derivative methods
described below share identical interpolation backends.

% -----------------------------------------------------------
\subsection{Derivative Methods for Bulk Chemical Potentials}
\label{sec:deriv_methods}

The phase-field model requires two thermodynamic quantities at each
grid point per time step: the bulk driving force
$\mu_j = \partial G_m/\partial c_j$ (first derivative) and the
Hessian $H_{ij} = \partial^2 G_m/\partial c_i\partial c_j$ (second
derivative, needed for stability analysis and time-step selection).
Three strategies are implemented for obtaining these quantities,
each representing a different trade-off between accuracy, data
requirements, and computational overhead.

\subsubsection{Method~I: Finite Differences on the $G_m$ Surrogate (\texttt{fd})}
\label{sec:fd}

The simplest approach is to differentiate the $G_m$ surrogate
numerically.  Central finite differences give
\begin{align}
  \mu_j^{\rm fd}
  &= \frac{G_m(\bm{c}+\varepsilon\bm{e}_j)
           - G_m(\bm{c}-\varepsilon\bm{e}_j)}{2\varepsilon}
  + O(\varepsilon^2),
  \label{eq:fd_mu}
  \\
  H_{ij}^{\rm fd}
  &= \begin{cases}
       \dfrac{G_m(\bm{c}+\varepsilon\bm{e}_i)
              - 2G_m(\bm{c})
              + G_m(\bm{c}-\varepsilon\bm{e}_i)}{\varepsilon^2}
       + O(\varepsilon^2) & i = j,\\[8pt]
       \dfrac{G_m(\bm{c}+\varepsilon\bm{e}_i+\varepsilon\bm{e}_j)
              - G_m(\bm{c}+\varepsilon\bm{e}_i-\varepsilon\bm{e}_j)
              - G_m(\bm{c}-\varepsilon\bm{e}_i+\varepsilon\bm{e}_j)
              + G_m(\bm{c}-\varepsilon\bm{e}_i-\varepsilon\bm{e}_j)}{4\varepsilon^2}
       + O(\varepsilon^2) & i \neq j,
     \end{cases}
  \label{eq:fd_H}
\end{align}
where $\bm{e}_j$ is the unit vector in the $j$-th composition
direction and $\varepsilon = 10^{-4}$ mol fraction.
This method requires only $G_m$ data (no chemical potential storage)
but introduces truncation errors that compound with the smoothing
bias of the $G_m$ surrogate.  In particular, the second-order
accurate stencil in Eq.~\eqref{eq:fd_H} can accumulate errors of
order $\varepsilon^2 G_m'''/G_m$ near the simplex boundaries where
$G_m$ is steep; our validation on the Nb--V binary at $500\,^\circ$C
yields an RMSE of $442\,\text{kJ\,mol}^{-1}$ in $H$ relative to the
directly computed CALPHAD Hessian (Section~\ref{sec:validation}).
This error is dominated by the entropy divergence
$H \to +\infty$ as $x \to 0$ or $x \to 1$ and its propagation
into the FD stencil.

\subsubsection{Method~II: Direct Interpolation of Jansson Derivatives (\texttt{jansson})}
\label{sec:jansson}

Under the simplex constraint $x_N = 1 - \sum_{i=1}^{N-1} c_i$, the
bulk chemical potential of field $j$ that enters the Cahn-Hilliard
equation is the \emph{relative} chemical potential
\begin{equation}
  \mu_j = \frac{\partial G_m}{\partial c_j}\bigg|_{\rm simplex}
        = \bar{\mu}_j - \bar{\mu}_N,
  \label{eq:mu_rel}
\end{equation}
where $\bar{\mu}_i = \texttt{MU}(i)$ is the absolute chemical
potential of element $i$~\cite{deFontaine1979}.  This identity
follows directly from the Gibbs-Duhem equation and is exact within
the single-phase thermodynamic model.

Similarly, the Hessian element is
\begin{equation}
  H_{ij}
  = \frac{\partial^2 G_m}{\partial c_i \partial c_j}\bigg|_{\rm simplex}
  = \frac{\partial(\bar{\mu}_i - \bar{\mu}_N)}{\partial x_j}
  = \Lambda_{ij} - \Lambda_{Nj},
  \label{eq:H_jansson}
\end{equation}
where $\Lambda_{ij} = \partial\bar{\mu}_i/\partial x_j$ is the
Jansson second derivative~\cite{jansson1984} computed by pycalphad
and stored in the HDF5 table (Section~\ref{sec:calphad}).

In the \texttt{jansson} mode, separate smooth surrogate functions are
built for each $\mu_j$ (Eq.~\ref{eq:mu_rel}) and each $H_{ij}$
(Eq.~\ref{eq:H_jansson}) using the same backend as the $G_m$
interpolant.  During time integration, all thermodynamic quantities
are obtained from these surrogates by simple function evaluation with
no numerical differentiation.  Because the Jansson derivatives are
computed by pycalphad using the analytic tangent to the Gibbs energy
surface, this method is free of FD truncation error and satisfies the
Gibbs-Duhem relation exactly at every composition point.  Our
validation confirms RMSE $= 0$ for both $\mu$ and $H$ relative to
the raw CALPHAD data (Section~\ref{sec:validation}).

For an $N$-component system, the \texttt{jansson} mode builds
$(N-1)$ surrogates for $\mu_j$ and $(N-1)^2$ surrogates for
$H_{ij}$, for a total of $(N-1)N$ additional interpolants
beyond the $G_m$ surface.  For quaternary alloys ($N=4$) this
amounts to 3 + 9 = 12 additional RBF$\to$RGI surfaces, each
constructed once at simulation startup.

\subsubsection{Method~III: Hybrid — Jansson $\mu$, Finite-Difference $H$ (\texttt{hybrid})}
\label{sec:hybrid}

The \texttt{hybrid} mode combines the best attributes of the two
preceding methods.  The bulk driving force $\mu_j$ is obtained by
direct interpolation of the stored Jansson $\bar{\mu}$ data
(Eq.~\ref{eq:mu_rel}), which eliminates the first-derivative
truncation error from the time-evolution equation.  The Hessian
$H_{ij}$ is obtained by finite differences on the $G_m$ surrogate
(Eq.~\ref{eq:fd_H}), which is used only for stability analysis and
time-step estimation rather than for the per-step driving force.
This is appropriate when the \texttt{dMU} dataset is absent or when
the raw Jansson data is suspected to be noisy (e.g., near a
miscibility gap or phase boundary where the single-phase free energy
surface becomes nearly flat and pycalphad's numerical derivatives are
less reliable).

A practical summary of the three methods is given in
Table~\ref{tab:methods}.

\begin{table}[htbp]
  \centering
  \caption{%
    Summary of derivative methods for bulk chemical potentials and
    Hessian in the CALPHAD-coupled Cahn-Hilliard model.
    RMSE values are for the Nb--V binary BCC$_{\rm A2}$ system at
    $T = 500\,^\circ$C with $n = 200$ composition grid points.
  }
  \label{tab:methods}
  \setlength{\tabcolsep}{6pt}
  \renewcommand{\arraystretch}{1.25}
  \begin{tabular}{llllcc}
    \hline
    Mode & $\mu_j$ source & $H_{ij}$ source
         & HDF5 data needed
         & RMSE$_\mu$ & RMSE$_H$ \\
    \hline
    \texttt{fd}
      & FD on $G_m$ (Eq.~\ref{eq:fd_mu})
      & FD on $G_m$ (Eq.~\ref{eq:fd_H})
      & \texttt{/GM}
      & 302 & 442\,200 \\
    \texttt{jansson}
      & Eq.~\ref{eq:mu_rel}
      & Eq.~\ref{eq:H_jansson}
      & \texttt{/GM}, \texttt{/MU}, \texttt{/dMU}
      & \textbf{0} & \textbf{0} \\
    \texttt{hybrid}
      & Eq.~\ref{eq:mu_rel}
      & FD on $G_m$ (Eq.~\ref{eq:fd_H})
      & \texttt{/GM}, \texttt{/MU}
      & \textbf{0} & 442\,200 \\
    \hline
  \end{tabular}
  \par\smallskip
  \footnotesize RMSE in J~mol$^{-1}$; \textbf{bold} = exact match
  with CALPHAD data.
\end{table}

% -----------------------------------------------------------
\subsection{Validation of Derivative Methods}
\label{sec:validation}

Figure~\ref{fig:deriv_methods} compares the three modes against raw
CALPHAD data for the Nb--V BCC$_{\rm A2}$ system at
$T = 500\,^\circ$C as a function of the vanadium mole fraction
$x_{\rm V}$.  The molar Gibbs energy $G_m$ (panel~a) is identical
across all modes because it is always obtained from the same
$G_m$ interpolant.  The relative chemical potential
$\mu_{\rm Nb} - \mu_{\rm V}$ (panel~b) shows that \texttt{fd}
deviates visibly from the CALPHAD reference near
$x_{\rm V} \approx 0.2$--$0.3$, the composition range where the
$G_m$ spline has the steepest first derivative.  Both \texttt{jansson}
and \texttt{hybrid} reproduce the CALPHAD chemical potentials exactly
at all tabulated compositions.

The most significant differences appear in the Hessian
$\partial^2 G_m/\partial x_{\rm Nb}^2$ (panel~c).  The y-axis is
clipped to $\pm 150\,\text{kJ\,mol}^{-1}$ to expose the spinodal
region (shaded pink, $H < 0$); the ideal-entropy divergence
$H \to +\infty$ as $x_{\rm V} \to 0$ or $1$ is annotated at both
edges.  The \texttt{fd} curve deviates substantially from the
CALPHAD reference due to the compounding of spline smoothing bias
and FD truncation error, with an RMSE of $442\,\text{kJ\,mol}^{-1}$
(Table~\ref{tab:methods}).  Since the Hessian controls the spinodal
decomposition wavelength via $\lambda^* = 4\pi\sqrt{\kappa/|H|}$,
errors of this magnitude translate directly into inaccurate
predictions of the dominant microstructural length scale.  The
\texttt{jansson} mode eliminates this error entirely, at the cost of
pre-computing the Jansson derivative dataset (step~1 of the workflow,
which is performed once and stored on disk).

We therefore recommend \texttt{jansson} mode as the default for
production runs, and \texttt{hybrid} mode as a fall-back when
the \texttt{dMU} dataset is unavailable or when Jansson data quality
near a phase boundary is suspect.  The \texttt{fd} mode is retained
for backward compatibility and for rapid prototyping where
thermodynamic accuracy is secondary to workflow simplicity.

% -----------------------------------------------------------
\subsection{Spinodal Stability and Interface Length Scales}
\label{sec:length_scales}

The spinodal decomposition instability is determined by the minimum
eigenvalue $\lambda_{\min}$ of the Hessian $\bm{H}(\bm{c}_0)$
evaluated at the mean composition $\bm{c}_0$.  If
$\lambda_{\min} < 0$, the alloy is spinodally unstable and phase
separation proceeds spontaneously.  The dominant decomposition
wavelength and interface width are
\begin{equation}
  \lambda_i^* = 4\pi\sqrt{\frac{\kappa_i}{|\lambda_{\min}|}},
  \qquad
  \delta_i = \sqrt{\frac{\kappa_i}{|\lambda_{\min}|}},
  \qquad
  i = 1,\ldots,N-1,
  \label{eq:length_scales}
\end{equation}
and the shortest unstable wavelength (critical wavelength below which
all modes are stable) is
\begin{equation}
  \lambda_{c,i} = 2\pi\sqrt{\frac{2\kappa_i}{|\lambda_{\min}|}}.
  \label{eq:lambda_c}
\end{equation}
Proper resolution of the interface requires $\delta_i > 5\Delta x$;
the gradient energy coefficient $\kappa_i$ must be chosen
accordingly.  A numerical stability criterion for the semi-implicit
scheme (Eq.~\ref{eq:semi_implicit}) is
$M_{ii}|\lambda_{\min}|\Delta t/\Delta x^2 < 0.5$, which is
evaluated per-field and reported alongside $\Delta t_{\rm stable}$
(Eq.~\ref{eq:dt_stable}) in the simulation parameter summary written
to disk.
